% ==============================================================================
% sci_appendix.tex — SCI/SICE tutorial appendix: cheat sheets, review paths,
% troubleshooting (take-home material)
% SCI/SICE チュートリアル 付録: チートシートと復習パス（持ち帰り資料）
% ==============================================================================

% ------------------------------------------------------------------
\begin{frame}[fragile]{sf CLI チートシート (1/2)}
  \begin{table}
    \centering\scriptsize
    \begin{tabular}{l>{\raggedright\arraybackslash}p{75mm}}
      \hline
      \textbf{コマンド} & \textbf{機能} \\
      \hline
      \code{sf doctor}     & 環境診断（問題があればまず実行） \\
      \code{sf build/flash/monitor} & ビルド・書き込み・シリアルモニタ \\
      \code{sf telemetry}  & 50\,Hz テレメトリのライブ表示（\code{--web} でブラウザ） \\
      \code{sf cal gyro/accel/mag} & センサキャリブレーション \\
      \code{sf lesson}     & 実習コードの管理（\code{switch sci2026:N}/\code{build}/\code{flash}，一覧は \code{list --course sci2026}） \\
      \code{sf trim analyze} & ホバーログから姿勢トリムを算出 \\
      \code{sf takeoff} 等   & Tello 風ミニコマンド（\code{sf takeoff/land/hover/jump}，個別のトップレベルコマンド） \\
      \hline
    \end{tabular}
  \end{table}
\end{frame}

% ------------------------------------------------------------------
\begin{frame}[fragile]{sf CLI チートシート (2/2)}
  \begin{table}
    \centering\scriptsize
    \begin{tabular}{l>{\raggedright\arraybackslash}p{75mm}}
      \hline
      \textbf{コマンド} & \textbf{機能} \\
      \hline
      \code{sf log wifi} & WiFi 経由で 400\,Hz テレメトリ取得（\code{capture} は旧ファーム専用） \\
      \code{sf log viz/analyze}  & ログ可視化・解析 \\
      \code{sf sim run vpython/genesis} & シミュレータ起動（セッション5） \\
      \code{sf sils build/scenario/gui} & SILS ベンチ（セッション5） \\
      \code{sf sysid fit}        & 閉ループ P 制御ログからプラント同定（実習7） \\
      \code{sf sysid rate-fit/rate-tune} & 閉ループ同定（ETFE + フィット）と仕様ベース PID 設計（セッション4） \\
      \code{sf sysid noise}      & 静止センサログから Allan 分散でノイズ特性を推定 \\
      \code{sf params check}     & 物理パラメータ整合検査 \\
      \code{sf upgrade}          & 最新版を pull し環境を再同期 \\
      \hline
    \end{tabular}
  \end{table}

  {\footnotesize \code{sf --help} の順序を基本に，授業で頻出する \code{sf lesson} を前に出している。全コマンドは \code{lib/sfcli/commands/} に実装がある}
\end{frame}

% ------------------------------------------------------------------
\begin{frame}[fragile]{sf lesson と通常コマンドの対応\later}
  {\footnotesize \code{sf lesson} は実習を滞りなく進めるための近道で，正は通常コマンド。各コマンドは次を呼んでいるだけ}

  \vspace{0.3em}

  \begin{table}
    \centering\footnotesize
    \setlength{\tabcolsep}{4pt}
    \begin{tabular}{>{\raggedright\arraybackslash}p{46mm}>{\raggedright\arraybackslash}p{86mm}}
      \hline
      \textbf{\code{sf lesson} コマンド} & \textbf{実体} \\
      \hline
      \code{sf lesson switch sci2026:N \allowbreak [--solution]} & 課題 N の \code{student.cpp}（\code{solution.cpp}）を実習ファームウェアの \code{user\_code.cpp} にコピー \\
      \code{sf lesson edit} & その \code{user\_code.cpp} を開く \\
      \code{sf lesson build} & \code{sf build workshop} \\
      \code{sf lesson flash} & \code{sf flash workshop -m}（\code{--no-monitor} で \code{-m} 無し） \\
      \code{sf lesson monitor} & \code{sf monitor workshop} \\
      \code{sf lesson sils} & \code{sf sils build \allowbreak --target workshop} $\to$ \code{sf sils scenario \allowbreak simulator/\allowbreak sils/\allowbreak scenarios/\allowbreak workshop\_acro.scn \allowbreak --target workshop} \\
      \code{sf lesson sils \allowbreak --scenario step} & 同上でシナリオは \code{workshop\_acro\_step.scn} \\
      \hline
    \end{tabular}
  \end{table}
\end{frame}

% ------------------------------------------------------------------
\begin{frame}[fragile]{ws:: API チートシート (1/2)}
  \begin{table}
    \centering\scriptsize
    \begin{tabular}{l>{\raggedright\arraybackslash}p{72mm}}
      \hline
      \textbf{カテゴリ} & \textbf{関数} \\
      \hline
      モータ制御 & \code{motor\_set\_duty(id,duty)} [id 1--4, duty 0--1], \code{motor\_set\_all(duty)} [0--1],
        \code{motor\_stop\_all()}, \code{motor\_mixer(t,r,p,y)} [t 0--1, r/p/y $-1$--1], \code{arm()},
        \code{disarm()}, \code{is\_armed()} \\
      コントローラ入力 & \code{rc\_throttle()} [0--1], \code{rc\_roll/pitch/yaw()} [$-1$--1] \\
      ボタン・モード & \code{rc\_throttle\_yaw\_button()}, \code{rc\_roll\_pitch\_button()},
        \code{rc\_stabilize\_acro\_mode()}, \code{rc\_alt\_mode()}, \code{rc\_pos\_mode()}\mbox{（すべて bool）} \\
      LED 制御 & \code{led\_color(r,g,b)} [0--255], \code{disable/enable\_led\_task()},
        \code{is\_led\_task\_disabled()} \\
      \hline
    \end{tabular}
  \end{table}

  {\footnotesize 実習中は Duty $\le$ 0.15 で運用する（安全上の運用ルール。ファーム側では強制されない）}
\end{frame}

% ------------------------------------------------------------------
\begin{frame}[fragile]{ws:: API チートシート (2/2)}
  \begin{table}
    \centering\scriptsize
    \begin{tabular}{l>{\raggedright\arraybackslash}p{72mm}}
      \hline
      \textbf{カテゴリ} & \textbf{関数} \\
      \hline
      IMU センサ & \code{gyro\_x/y/z()} [rad/s], \code{accel\_x/y/z()} [m/s$^2$] \\
      環境・距離センサ & \code{baro\_altitude/pressure()} [m, Pa], \code{mag\_x/y/z()} [$\mu$T],
        \code{tof\_bottom/front()} [m], \code{flow\_vx/vy/quality()} [m/s, 0--255] \\
      姿勢推定 & \code{estimated\_roll/pitch/yaw()} [rad], \code{estimated\_altitude()} [m] \\
      制御目標のロギング & \code{set\_rate\_target(r,p,y)} [rad/s], \code{set\_angle\_target(r,p)} [rad]
        （実習7，Data Stream の \code{rate\_ref\_*}/\code{angle\_ref\_*} に記録，制御には無関係） \\
      ユーティリティ & \code{millis()} [ms], \code{battery\_voltage()} [V], \code{print(fmt,...)},
        \code{set\_channel(ch)} [1, 6, 11] \\
      \hline
    \end{tabular}
  \end{table}

  {\footnotesize 実習用 API のヘッダ \code{\#include "workshop\_api.hpp"}，全関数は \code{ws::} 名前空間}
\end{frame}

% ------------------------------------------------------------------
\begin{frame}{トラブルシューティング (1/2)}
  \begin{table}
    \centering\scriptsize
    \begin{tabular}{>{\raggedright\arraybackslash}p{28mm}>{\raggedright\arraybackslash}p{28mm}>{\raggedright\arraybackslash}p{55mm}}
      \hline
      \textbf{症状} & \textbf{原因} & \textbf{対処} \\
      \hline
      StampFly の WiFi が見つからない & 電源が入っていない & バッテリーを確認，電源を入れ直す \\
      接続しても通信できない & IP アドレスが違う & \code{192.168.10.1} を確認 \\
      離陸しない & キャリブレーション未実施 & \code{sf cal gyro} を実行 \\
      振動する & PID ゲインが高すぎる & ゲインを下げる（セッション4参照） \\
      シミュレータに自動フォールバック & WiFi 未接続 & 意図的ならそのまま使用可 \\
      \hline
    \end{tabular}
  \end{table}

  {\footnotesize 出典: \code{docs/guides/troubleshooting.md}}
\end{frame}

% ------------------------------------------------------------------
\begin{frame}{トラブルシューティング (2/2)}
  \begin{table}
    \centering\scriptsize
    \begin{tabular}{>{\raggedright\arraybackslash}p{28mm}>{\raggedright\arraybackslash}p{28mm}>{\raggedright\arraybackslash}p{55mm}}
      \hline
      \textbf{症状} & \textbf{原因} & \textbf{対処} \\
      \hline
      ビルドエラー（赤字 \code{error:}） & 打ち間違い（カンマ・セミコロン抜け等） & エラー1行目のファイル名・行番号を確認し実習コードと見比べる \\
      \code{sf lesson flash} が失敗 & 書き込みポートの競合・接触不良 & 別のケーブル／ポートを試す。改善しなければ機体の BOOT ボタンを押したまま USB を挿し直し，書き込み専用モードで再実行する \\
      モータが回らない & ARM されていない，または USB 給電中（安全のため ARM 拒否） & 機体ボタンをクリックして ARM 状態を確認する。USB を外してバッテリー駆動になっているか確認する \\
      起動 LED が緑にならない & 静止校正中に機体を動かしてしまった & 機体を静かに置き直し，緑点灯とビープ3回（readyTone）を待つ \\
      \hline
    \end{tabular}
  \end{table}
\end{frame}

% ------------------------------------------------------------------
% NOTE: 元は1枚・10行の表だったが、footnotesize では frame の縦幅（vbox）に
% 収まらず（実測で規定の縦幅を約13mm超過）、scriptsizeに落としても\code
% マクロが常に\smallで組まれるため解消しなかった。列幅・行間・余白の調整
% だけでは収まらないため、sf CLIチートシート等と同じ(1/2)/(2/2)の2枚構成に
% 分割し、footnotesize を維持する。
\begin{frame}{5セッション復習パス (1/2)}
  \begin{table}
    \centering\footnotesize
    \setlength{\tabcolsep}{4pt}
    \begin{tabular}{l>{\raggedright\arraybackslash}p{40mm}>{\raggedright\arraybackslash}p{33mm}>{\raggedright\arraybackslash}p{55mm}}
      \hline
      \textbf{セッション} & \textbf{コマンド} & \textbf{扱う内容} & \textbf{文書} \\
      \hline
      1 & \code{sf sim run vpython}        & シミュレータ操縦        & \code{docs/\allowbreak commands/\allowbreak sf-sim.md} \\
      1 & \code{sf sils gui}               & SILS の体験             & \code{simulator/\allowbreak sils/\allowbreak gui/\allowbreak README.md} \\
      2 & \code{sf lesson switch sci2026:2} & IMU センサ（実習2）    & \code{firmware/\allowbreak vehicle/\allowbreak docs/\allowbreak architecture.md \S2} \\
      2 & \code{sf log wifi/viz}           & テレメトリ取得と可視化  & \code{docs/\allowbreak guides/\allowbreak flight-log-viz.md} \\
      3 & \code{sf lesson switch sci2026:3} & モータ制御（実習3）     & \code{firmware/\allowbreak vehicle/\allowbreak docs/\allowbreak hardware\_init.md} \\
      3 & \code{sf lesson switch sci2026:4} & コントローラ入力（実習4） & \code{protocol/\allowbreak spec/}（ControlPacket） \\
      \hline
    \end{tabular}
  \end{table}
\end{frame}

% ------------------------------------------------------------------
\begin{frame}{5セッション復習パス (2/2)}
  \begin{table}
    \centering\footnotesize
    \setlength{\tabcolsep}{4pt}
    \begin{tabular}{l>{\raggedright\arraybackslash}p{40mm}>{\raggedright\arraybackslash}p{33mm}>{\raggedright\arraybackslash}p{55mm}}
      \hline
      \textbf{セッション} & \textbf{コマンド} & \textbf{扱う内容} & \textbf{文書} \\
      \hline
      4 & \code{sf lesson switch sci2026:8} & PID 制御（実習8）       & 本資料セッション4「PID 制御器の構成」〜「ARM 遷移時の積分リセット」 \\
      4 & \code{sf sysid fit}              & システム同定（実習7）    & 本資料セッション4「システム同定の考え方」〜「同定結果と理論値の比較」 \\
      5 & \code{sf sils gui}               & SILS ブラウザ実験環境   & \code{simulator/\allowbreak sils/\allowbreak gui/\allowbreak README.md} \\
      5 & \code{sf log viz}                & ログ解析                & \code{docs/\allowbreak guides/\allowbreak flight-log-viz.md} \\
      \hline
    \end{tabular}
  \end{table}
\end{frame}

% ------------------------------------------------------------------
\begin{frame}{持ち帰りのご案内}
  \begin{exampleblock}{資料はすべて公開されている}\footnotesize
    \begin{itemize}
      \item リポジトリ: \url{https://github.com/M5Fly-kanazawa/stampfly_ecosystem}
      \item ドキュメント: \url{https://www.docswell.com/s/Kouhei_Ito/5L387D-2026-09-10-064635}
      \item 本日の内容はオンデマンド配信でも後日視聴できる
    \end{itemize}
  \end{exampleblock}

  \vspace{0.3em}

  \begin{block}{ご自身の StampFly で}\footnotesize
    ノート PC と StampFly さえあれば，今日のデモは自宅・研究室でそのまま再現できる。
    まずは \code{sf doctor} で環境を確認してから始めてほしい。
    ご質問・不具合報告はリポジトリの Issue へ（開発協力も歓迎）
  \end{block}
\end{frame}
